\section{Ejercicio 2}
\newcommand{\ft}[1]{\color[HTML]{CBCEFB} #1}

\subsection{UxU}

\subsubsection{Enteros}
\begin{table}[H]
    \centering
    \begin{tabular}{|l|c|c|c|c|c|c|c|c|c|c|c|c|c|c|c|c|r|}
                                                                                                                                    \hline
        U(8, 0)  &        &        &        &        &        &        &        &        & 1 & 0 & 1 & 1 & 0 & 0 & 1 & 1 &   179 \\ \hline
        U(8, 0)  &        &        &        &        &        &        &        &        & 1 & 1 & 0 & 0 & 0 & 1 & 0 & 1 &   197 \\ \hline
                 & \ft{0} & \ft{0} & \ft{0} & \ft{0} & \ft{0} & \ft{0} & \ft{0} & \ft{0} & 1 & 0 & 1 & 1 & 0 & 0 & 1 & 1 &       \\ \cline{2-17}
                 & \ft{0} & \ft{0} & \ft{0} & \ft{0} & \ft{0} & \ft{0} & \ft{0} & 0      & 0 & 0 & 0 & 0 & 0 & 0 & 0 & - &       \\ \cline{2-17}
                 & \ft{0} & \ft{0} & \ft{0} & \ft{0} & \ft{0} & \ft{0} &      1 & 0      & 1 & 1 & 0 & 0 & 1 & 1 & - & - &       \\ \cline{2-17}
                 & \ft{0} & \ft{0} & \ft{0} & \ft{0} & \ft{0} & 0      &      0 & 0      & 0 & 0 & 0 & 0 & 0 & - & - & - &       \\ \cline{2-17}
                 & \ft{0} & \ft{0} & \ft{0} & \ft{0} &      0 & 0      &      0 & 0      & 0 & 0 & 0 & 0 & - & - & - & - &       \\ \cline{2-17}
                 & \ft{0} & \ft{0} & \ft{0} & 0      &      0 & 0      &      0 & 0      & 0 & 0 & 0 & - & - & - & - & - &       \\ \cline{2-17}
                 & \ft{0} & \ft{0} &      1 & 0      &      1 & 1      &      0 & 0      & 1 & 1 & - & - & - & - & - & - &       \\ \cline{2-17}
                 & \ft{0} & 1      &      0 & 1      &      1 & 0      &      0 & 1      & 1 & - & - & - & - & - & - & - &       \\ \hline
        U(16, 0) &      1 & 0      &      0 & 0      &      1 & 0      &      0 & 1      & 1 & 0 & 1 & 1 & 1 & 1 & 1 & 1 & 35263 \\ \hline
    \end{tabular}
\end{table}

\subsubsection{Ambos U(8,7)}
Al resultado anterior, se dividen por $2^{14}$ (7 bits fraccionales por elemento) y se obtiene: $10.00100110111111$, que es $2.15228271484375$ en base 10. 

\subsubsection{A entero, B S(8,7)}

\subsection{UxS}

\subsubsection{Enteros}
\begin{table}[H]
    \centering
    \begin{tabular}{|l|c|c|c|c|c|c|c|c|c|c|c|c|c|c|c|c|r|}
                                                                                                                                    \hline
        U(8, 0) &        &        &        &        &        &        &        &        & 1 & 0 & 1 & 1 & 0 & 0 & 1 & 1 &    179 \\ \hline
        S(8, 0) &        &        &        &        &        &        &        &        & 1 & 1 & 0 & 0 & 0 & 1 & 0 & 1 &    -59 \\ \hline
                & \ft{0} & \ft{0} & \ft{0} & \ft{0} & \ft{0} & \ft{0} & \ft{0} & \ft{0} & 1 & 0 & 1 & 1 & 0 & 0 & 1 & 1 &        \\ \cline{2-17}
                & \ft{0} & \ft{0} & \ft{0} & \ft{0} & \ft{0} & \ft{0} & \ft{0} & 0      & 0 & 0 & 0 & 0 & 0 & 0 & 0 & - &        \\ \cline{2-17}
                & \ft{0} & \ft{0} & \ft{0} & \ft{0} & \ft{0} & \ft{0} &      1 & 0      & 1 & 1 & 0 & 0 & 1 & 1 & - & - &        \\ \cline{2-17}
                & \ft{0} & \ft{0} & \ft{0} & \ft{0} & \ft{0} & 0      &      0 & 0      & 0 & 0 & 0 & 0 & 0 & - & - & - &        \\ \cline{2-17}
                & \ft{0} & \ft{0} & \ft{0} & \ft{0} &      0 & 0      &      0 & 0      & 0 & 0 & 0 & 0 & - & - & - & - &        \\ \cline{2-17}
                & \ft{0} & \ft{0} & \ft{0} & 0      &      0 & 0      &      0 & 0      & 0 & 0 & 0 & - & - & - & - & - &        \\ \cline{2-17}
                & \ft{0} & \ft{0} &      1 & 0      &      1 & 1      &      0 & 0      & 1 & 1 & - & - & - & - & - & - &        \\ \cline{2-17}
                & \ft{1} & 0      &      1 & 0      &      0 & 1      &      1 & 0      & 1 & - & - & - & - & - & - & - &        \\ \hline
        S(16,0) &      1 & 1      &      0 & 1      &      0 & 1      &      1 & 0      & 1 & 0 & 1 & 1 & 1 & 1 & 1 & 1 & -10561 \\ \hline
    \end{tabular}
\end{table}

\subsubsection{Ambos S(8,7)}
Al resultado anterior, se dividen por $2^{14}$ (7 bits fraccionales por elemento) y se obtiene: $11.01011010111111$, que es $-0.64459228515625$ en base 10. 

\subsubsection{A entero, B S(8,7)}

\subsection{SxU}

\subsubsection{Enteros}
\begin{table}[H]
    \centering
    \begin{tabular}{|l|c|c|c|c|c|c|c|c|c|c|c|c|c|c|c|c|r|}
                                                                                                                                     \hline
        S(8, 0)  &        &        &        &        &        &        &        &        & 1 & 0 & 1 & 1 & 0 & 0 & 1 & 1 &    -77 \\ \hline
        U(8, 0)  &        &        &        &        &        &        &        &        & 1 & 1 & 0 & 0 & 0 & 1 & 0 & 1 &    197 \\ \hline
                 & \ft{1} & \ft{1} & \ft{1} & \ft{1} & \ft{1} & \ft{1} & \ft{1} & \ft{1} & 1 & 0 & 1 & 1 & 0 & 0 & 1 & 1 &        \\ \cline{2-17}
                 & \ft{0} & \ft{0} & \ft{0} & \ft{0} & \ft{0} & \ft{0} & \ft{0} & 0      & 0 & 0 & 0 & 0 & 0 & 0 & 0 & - &        \\ \cline{2-17}
                 & \ft{1} & \ft{1} & \ft{1} & \ft{1} & \ft{1} & \ft{1} &      1 & 0      & 1 & 1 & 0 & 0 & 1 & 1 & - & - &        \\ \cline{2-17}
                 & \ft{0} & \ft{0} & \ft{0} & \ft{0} & \ft{0} & 0      &      0 & 0      & 0 & 0 & 0 & 0 & 0 & - & - & - &        \\ \cline{2-17}
                 & \ft{0} & \ft{0} & \ft{0} & \ft{0} &      0 & 0      &      0 & 0      & 0 & 0 & 0 & 0 & - & - & - & - &        \\ \cline{2-17}
                 & \ft{0} & \ft{0} & \ft{0} & 0      &      0 & 0      &      0 & 0      & 0 & 0 & 0 & - & - & - & - & - &        \\ \cline{2-17}
                 & \ft{1} & \ft{1} &      1 & 0      &      1 & 1      &      0 & 0      & 1 & 1 & - & - & - & - & - & - &        \\ \cline{2-17}
                 & \ft{1} & 1      &      0 & 1      &      1 & 0      &      0 & 1      & 1 & - & - & - & - & - & - & - &        \\ \hline
        S(16, 0) &      1 & 1      &      0 & 0      &      0 & 1      &      0 & 0      & 1 & 0 & 1 & 1 & 1 & 1 & 1 & 1 & -15169 \\ \hline
    \end{tabular}
\end{table}

\subsubsection{Ambos S(8,7)}
Al resultado anterior, se dividen por $2^{14}$ (7 bits fraccionales por elemento) y se obtiene: $11.00010010111111$, que es $-0.92584228515625$ en base 10. 

\subsubsection{A entero, B S(8,7)}

\subsection{SxS}

\subsubsection{Enteros}
\begin{table}[H]
    \centering
    \begin{tabular}{|l|c|c|c|c|c|c|c|c|c|c|c|c|c|c|c|c|r|}
                                                                                                                                  \hline
        S(8, 0) &        &        &        &        &        &        &        &        & 1 & 0 & 1 & 1 & 0 & 0 & 1 & 1 &  -77 \\ \hline
        S(8, 0) &        &        &        &        &        &        &        &        & 1 & 1 & 0 & 0 & 0 & 1 & 0 & 1 &  -59 \\ \hline
                & \ft{1} & \ft{1} & \ft{1} & \ft{1} & \ft{1} & \ft{1} & \ft{1} & \ft{1} & 1 & 0 & 1 & 1 & 0 & 0 & 1 & 1 &      \\ \cline{2-17}
                & \ft{0} & \ft{0} & \ft{0} & \ft{0} & \ft{0} & \ft{0} & \ft{0} & 0      & 0 & 0 & 0 & 0 & 0 & 0 & 0 & - &      \\ \cline{2-17}
                & \ft{1} & \ft{1} & \ft{1} & \ft{1} & \ft{1} & \ft{1} &      1 & 0      & 1 & 1 & 0 & 0 & 1 & 1 & - & - &      \\ \cline{2-17}
                & \ft{0} & \ft{0} & \ft{0} & \ft{0} & \ft{0} & 0      &      0 & 0      & 0 & 0 & 0 & 0 & 0 & - & - & - &      \\ \cline{2-17}
                & \ft{0} & \ft{0} & \ft{0} & \ft{0} &      0 & 0      &      0 & 0      & 0 & 0 & 0 & 0 & - & - & - & - &      \\ \cline{2-17}
                & \ft{0} & \ft{0} & \ft{0} & 0      &      0 & 0      &      0 & 0      & 0 & 0 & 0 & - & - & - & - & - &      \\ \cline{2-17}
                & \ft{1} & \ft{1} &      1 & 0      &      1 & 1      &      0 & 0      & 1 & 1 & - & - & - & - & - & - &      \\ \cline{2-17}
                & \ft{0} & 0      &      1 & 0      &      0 & 1      &      1 & 0      & 1 & - & - & - & - & - & - & - &      \\ \hline
        S(16,0) &      0 & 0      &      0 & 1      &      0 & 0      &      0 & 1      & 1 & 0 & 1 & 1 & 1 & 1 & 1 & 1 & 4543 \\ \hline
    \end{tabular}
\end{table}

\subsubsection{Ambos S(8,7)}
Al resultado anterior, se dividen por $2^{14}$ (7 bits fraccionales por elemento) y se obtiene: $00.01000110111111$, que es $0.27728271484375$ en base 10. 

\subsubsection{A entero, B S(8,7)}

